Several limitations were encountered and are reported honestly here.

\subsection{Low-Mach Dissipation of the Rusanov Flux}
The dominant accuracy limitation of the production scheme is the
$O(a)$ dissipation of the Rusanov flux at low Mach number. At $M=0.1$--0.15
the artificial dissipation $s_{\max}\,\Delta\mathbf U$ overwhelms the physical
viscous terms in the boundary layer (an effective-Reynolds-number reduction).
Consequences visible in the results:
\begin{itemize}
\item Laminar skin friction on the NACA0012 cases is over-predicted by
roughly an order of magnitude relative to Blasius-level expectations
(e.g.\ \texttt{naca0012\_m015\_laminar\_re5000} $C_D=0.49$), and the
\texttt{cylinder\_m010\_laminar\_re20} drag ($C_D\approx6.0$) is about
$3\times$ the literature value of $\approx2.0$ for a steady Re\,20 cylinder
wake. Pressure-dominated quantities (inviscid and supersonic wave drag, lift
symmetry) are accurate; viscous-dominated drags at low Mach are qualitative
only.
\item The same over-dissipation suppresses the Re\,200 vortex street when
the Rusanov flux is used in the transient: the wake is stable and the drag
settles to an unphysical $\approx4.6$. The production Re\,200 run therefore
uses the Roe flux with a Harten--Yee entropy fix in the transient driver
only (see \cref{sec:re200}), which resolves shedding with physical
amplitudes.
\end{itemize}
Two remedies were implemented and evaluated. (a) A low-Mach dissipation
scaling $\chi=\min(1,\max(M_L,M_R)/M_{ref})$ with $M_{ref}=0.3$ applied to
the acoustic part of $s_{\max}$ improves low-Mach drag substantially
(Re\,200 shedding appears with $C_L\approx\pm0.5$, laminar drags drop toward
literature values) but destabilizes the steady implicit runs: the reduced
dissipation interacts with the LU-SGS scalar Jacobian, producing a stall or
slow divergence in the steady NACA cases. (b) The Roe flux was tried for the
steady cases and diverges from a freestream start on these meshes regardless
of entropy-fix strength. Both experiments are documented in the repository
history. The production steady cases therefore keep the full-strength
Rusanov dissipation ($\chi=1$), accepting the viscous-drag error, and the
limitation is confined to low-Mach viscous-drag accuracy.

\subsection{Residual Limit Cycles in Transonic Cases}
The transonic inviscid case (\texttt{naca0012\_m080\_inviscid}) and, at some
partition counts, the supersonic cases settle into a bounded residual limit
cycle (L2 residual oscillating around $2\times10^{-6}$, about 2 orders of
reduction) instead of converging to the 3-order target. This is the classic
limiter--shock interaction: the Venkatakrishnan limiter and the shock
displacement exchange energy each pseudo-step. Force coefficients are
stationary to four digits inside the cycle, the Mach field is clean, and the
run-status notes record the plateau stop. The affected runs are marked
\texttt{converged} on the basis of stable forces, which the metadata reports
together with the achieved reduction orders.

\subsection{Sensitivity to Partition Count}
Because LU-SGS relaxation is block-Jacobi across ranks (ghost increments lag
one sweep), the nonlinear convergence path depends mildly on the partition:
the same case can converge directly at $np=4$ but land in the limit cycle at
$np=8$, changing the step count (though not the converged forces, see
\cref{tab:rankstudy}). This is a known property of domain-decomposed
implicit smoothers and is reported rather than hidden.

\subsection{Other Limitations}
\begin{itemize}
\item The update clipping (10\% maximum relative change in $\rho$ and $p$
per step) is a robustness device for the impulsive freestream start; it is
inactive once the residual has dropped by about an order.
\item The farfield of the supplied NACA mesh is only about 2 chords away,
so the inviscid drags include a small farfield-boundary effect; no attempt
was made to correct for it.
\item Restart files are per-rank binary snapshots tied to the partition;
restarting at a different rank count requires repartitioning and is not
supported.
\end{itemize}
